% TeX 风格 DSL - 语法参考
% LaTeX 格式的 Galgame 脚本

% ============================================
% 语法规则说明
% ============================================

/*
1. 角色定义：\character[ID, 显示名][简写]
2. 变量定义：\variable(变量名 = 初始值)
3. 场景定义：\scene[背景图, 音乐]
4. 角色显示：\char[角色ID][表情, 位置]
5. 对话显示：\简写[表情] 对话内容
6. 选择分支：\choice ... \case[目标][效果] 文本 ... \end-choice
7. 条件选项：\case[目标](条件) 文本
8. 变量操作：\set(变量名 = 值)
9. 跳转：\goto(目标)
10. 调用：\load(文件名)
11. 注释：% 单行注释 或 /* 多行注释 */
*/

% ============================================
% 角色定义
% ============================================

\character[player, 主角][p]
\character[sakura, 樱井美咲][s]
\character[ren, 黑崎莲][r]
\character[yuki, 白雪悠希][y]

% ============================================
% 变量定义
% ============================================

\variable(sakura_affection = 0)
\variable(ren_affection = 0)
\variable(yuki_affection = 0)
\variable(trust_points = 0)
\variable(clues_found = 0)
\variable(knows_secret = false)
\variable(knows_ren_secret = false)
\variable(has_lab_key = false)
\variable(has_lab_access = false)
\variable(school_route = false)
\variable(home_route = false)
\variable(ending_type = "")

% ============================================
% 完整示例场景
% ============================================

\label(example_scene)

\scene[classroom.png, peaceful.mp3]

\char[sakura][smile, center]

\s 你好！欢迎来到星光学院！

\choice
\case[next_part][+s] 你好！
\case[introduce] 请问你是？
\case[secret_talk](knows_secret) 关于那个秘密...
\end-choice

% ============================================
% 分支：你好
% ============================================

\label(next_part)

\p 你好！

\s[smile] 要不要参观一下学校？

\choice
\case[tour] 好啊
\case[end] 下次吧
\end-choice

% ============================================
% 分支：介绍
% ============================================

\label(introduce)

\s 我是樱井美咲，学生会副会长！

\goto(next_part)

% ============================================
% 分支：秘密
% ============================================

\label(secret_talk)

\p 关于那个秘密...

\goto(next_part)

% ============================================
% 分支：参观
% ============================================

\label(tour)

\scene[school_hallway.png]

\s 这里是教学楼，那边是图书馆...

\goto(end)

% ============================================
% 结束
% ============================================

\label(end)

\s 再见！

\fadeout(1.0)

\return

% ============================================
% 复杂分支示例
% ============================================

\label(complex_choice)

\r 你想知道真相吗？

\choice
\case[reveal_all][+t, +r] 我想知道一切
\case[avoid_danger][-t] 太危险了，算了吧
\case[share_intel](clues_found >= 3)[+r] 我已经知道一些了
\case[ren_confession](knows_ren_secret)[+r] 你也是实验对象？
\end-choice

% ============================================
% 多条件判断示例
% ============================================

\label(determine_ending)

\if(sakura_affection >= 8 && trust_points >= 10)
\goto(sakura_true_ending)
\elif(ren_affection >= 8 && knows_ren_secret)
\goto(ren_true_ending)
\elif(yuki_affection >= 6 && clues_found >= 5)
\goto(yuki_true_ending)
\elif(sakura_affection >= 5)
\goto(sakura_good_ending)
\elif(ren_affection >= 5)
\goto(ren_good_ending)
\elif(yuki_affection >= 5)
\goto(yuki_good_ending)
\elif(trust_points >= 8)
\goto(good_ending)
\else
\goto(normal_ending)
\endif

% ============================================
% 结局定义
% ============================================

\label(sakura_good_ending)

\scene[cherry_blossom.png, romantic.mp3]

\char[sakura][blush, center]

\s 那个...我有话想对你说。

\p 什么事？

\s[shy] 这段时间...谢谢你一直陪在我身边。

\s 我...我好像...

\p ？

\s[determined] 我喜欢你！

\effect(heart_float, 1.0)

\p 樱井同学...

\s[happy] 以后...请多多指教了！

\fadeout(2.0)

\ending(樱井美咲结局：樱花之约)

你与樱井美咲成为了恋人。
在樱花树下，你们许下了永远的约定。

\goto(show_stats)

% ============================================
% 游戏统计
% ============================================

\label(show_stats)

---

游戏统计

* 樱井美咲好感度: \print(sakura_affection)
* 黑崎莲好感度: \print(ren_affection)
* 白雪悠希好感度: \print(yuki_affection)
* 信任点数: \print(trust_points)
* 发现线索: \print(clues_found)

---

感谢游玩《星光学院的秘密》！

\return

% ============================================
% 常用命令参考
% ============================================

% 音频命令
% \scene[背景, 音乐] - 设置场景和音乐
% \audio.play(文件路径, 音量) - 播放音频
% \audio.stop() - 停止音频
% \audio.fadeout(时长) - 淡出

% 场景命令
% \scene[背景图] - 切换背景
% \transition(效果, 时长) - 过渡效果

% 角色命令
% \char[角色][表情, 位置] - 显示角色
% \char[角色][hide] - 隐藏角色
% 位置：left, center, right

% 特效命令
% \effect(效果名, 时长) - 播放特效
% \flash(时长) - 闪光
% \shake(时长, 强度) - 震动
% \fadeout(时长) - 淡出

% 变量操作
% \set(变量 = 值) - 设置变量
% \set(变量 += 值) - 增加变量
% \print(变量) - 输出变量值

% 流程控制
% \label(名称) - 定义标签
% \goto(目标) - 跳转
% \load(文件) - 加载文件
% \return - 返回
% \if(条件) ... \elif(条件) ... \else ... \endif - 条件判断
